% Pista alg-26j-q2 · Àlgebra
% Procedència: PAU Catalunya 2026 · Convocatòria de juny · Sèrie 1 · Exercici 2
\documentclass[11pt,a4paper]{article}
\usepackage{pista}
\pistainfo{Catalunya 2026}{Convocatòria de juny}{Sèrie 1}

\begin{document}

\section*{\textcolor{blauPista}{Exercici 2}}

\noindent Considereu el sistema d'equacions lineals següent, que està format per tres plans a l'espai i depèn del paràmetre real $m$:
\[
\left.\begin{array}{r}
x + my + z = 4 \\
x + 3y + z = 5 \\
mx + y + z = 4
\end{array}\right\}.
\]

\subsection*{2.1. \normalfont Discutiu el sistema per als diferents valors del paràmetre $m$. \hspace*{\fill}\textnormal{[1~punt]}}

\pista{Discuteix el sistema amb el teorema de Rouché-Frobenius:
\begin{itemize}
\item escriu la matriu de coeficients $A$ i la matriu ampliada $A^{'}$;
\item calcula $\det(A)$ i troba els valors de $m$ que l'anul·len (t'ha de sortir una equació de segon grau en $m$, amb dues arrels);
\item per a cadascun d'aquests valors, compara els rangs d'$A$ i d'$A^{'}$ per decidir si el sistema és compatible indeterminat o incompatible;
\item no t'oblidis del cas $\det(A) \neq 0$: quin tipus de sistema és?.
\end{itemize}}

\subsection*{2.2. \normalfont Interpreteu geomètricament aquest sistema per a tots els valors del paràmetre $m$ i resoleu-lo, si és possible, per al cas $m = 1$. \hspace*{\fill}\textnormal{[1~punt]}}

\pista{Cada equació representa un pla de l'espai: estàs estudiant la posició relativa de tres plans.
\begin{itemize}
\item Per a cadascun dels tres tipus de sistema de l'apartat anterior, pensa què tenen en comú els tres plans: quants punts i, si n'hi ha infinits, quina figura formen.
\item En el cas incompatible, fixa't si, per algun valor, dos dels plans són paral·lels.
\end{itemize}
\textit{Resoldre per a $m = 1$.} Comprova primer a quin dels tres casos correspon. Si és compatible indeterminat, allibera una variable (per exemple $z = \lambda$) i expressa-hi la resta d'incògnites.}

\subsection*{2.3. \normalfont Per a $m = 1$, és possible afegir una quarta equació de manera que el sistema resultant sigui compatible determinat i tingui com a solució $(x, y, z) = \left(3, \dfrac{1}{2}, \dfrac{1}{2}\right)$? Raoneu la resposta. \hspace*{\fill}\textnormal{[0,5~punts]}}

\pista{Ho pots fer en tres passos:
\begin{itemize}
\item comprova que el punt $\bigl(3, \frac{1}{2}, \frac{1}{2}\bigr)$ és una de les infinites solucions de l'apartat anterior per a $m = 1$: l'has de poder obtenir amb un valor concret de $\lambda$;
\item pensa quines dues condicions ha de complir la quarta equació perquè aquest punt sigui l'\emph{única} solució: una té a veure amb el punt, i l'altra, amb el rang de la matriu de coeficients, que ara té quatre files;
\item proposa una equació ben senzilla (per exemple, fixant el valor d'una incògnita) i comprova que compleix les dues condicions.
\end{itemize}

\medskip
\textbf{Atenció:} no totes les incògnites serveixen per fixar-ne el valor.}

\end{document}
